\section{Problem Formulation}
\label{sec:problem}

At environment step $t$, a frozen action-chunk policy receives observation
$o_t$ and predicts $H$ future actions,
\begin{equation}
\hat{\mathbf A}_t=\pi_\theta(o_t)
=\bigl[\hat{\mathbf a}_{t\mid t},\ldots,
\hat{\mathbf a}_{t+H-1\mid t}\bigr],
\label{eq:chunk}
\end{equation}
where $\theta$ is fixed during all experiments.  Each action is partitioned into
$G$ physically meaningful groups,
$\hat{\mathbf a}_{\tau\mid t}=
[\hat{\mathbf a}^{(1)}_{\tau\mid t},\ldots,
\hat{\mathbf a}^{(G)}_{\tau\mid t}]$.
For the present LIBERO study, $G=2$: six relative arm/end-effector components
and one gripper component.

A global fixed execution horizon $h$ queries the policy, executes the first $h$
full actions of $\hat{\mathbf A}_t$, and then repeats from the new observation.
Thus all groups cross a chunk boundary at the same environment step.  The design
question is whether the groups may instead use
\begin{equation}
\mathbf h=[h_1,\ldots,h_G],
\label{eq:group_horizons}
\end{equation}
while the joint policy and observation stream remain unchanged.  We study this
question through fixed schedules only.  Learning or choosing $\mathbf h$ online
is outside the current method.

We assume that the groups are disjoint and cover the action vector, and that
$1\leq h_g\leq H$.  The partition is supplied by known controller semantics,
not inferred from rollouts.  A diagonal vector
$\mathbf h=[h,\ldots,h]$ is exactly the global rule: every group expires
together, accepts the same fresh chunk, and advances at the same local index.
Off-diagonal vectors change only the acceptance schedule; observation
preprocessing, policy inference, and low-level action application remain fixed.

More specifically, $h_g$ is the number of environment steps for which group
$g$ retains its currently accepted predicted subsequence before accepting a fresh
subsequence.  It is not an independent policy-query interval for group $g$.
Policy observation and querying remain shared: whenever any group expires, the
frozen policy is queried once for a new full joint chunk.  Thus, for example,
with $(h_a,h_g)=(4,16)$, arm expiry can trigger a full-policy query every four
environment steps while the gripper retains its accepted subsequence for sixteen
steps before replacing it.
